\documentclass[11pt,a4paper]{article}
\usepackage[T1]{fontenc}
\usepackage{lmodern}
\usepackage[margin=27mm]{geometry}
\usepackage{amsmath,amssymb,amsthm,mathtools}
\usepackage{microtype}
\usepackage{booktabs,array}
\usepackage{enumitem}
\usepackage[hidelinks]{hyperref}
\setlength{\parskip}{5pt plus 1pt}
\setlength{\parindent}{0pt}
\setlength{\emergencystretch}{3em}
\setlist{itemsep=2pt,topsep=3pt}
\raggedbottom
\newtheorem{theorem}{Theorem}[section]
\newtheorem{lemma}[theorem]{Lemma}
\newtheorem{corollary}[theorem]{Corollary}
\newtheorem{proposition}[theorem]{Proposition}
\theoremstyle{definition}
\newtheorem{definition}[theorem]{Definition}
\newtheorem{example}[theorem]{Example}
\theoremstyle{remark}
\newtheorem{remark}[theorem]{Remark}
\newcommand{\Z}{\mathbb{Z}}
\newcommand{\N}{\mathbb{N}}
\newcommand{\Q}{\mathbb{Q}}
\newcommand{\F}{\mathbb{F}}
\newcommand{\ps}[1]{#1[[x]]}
\newcommand{\coeff}[2]{[x^{#1}]\,#2}
\newcommand{\hl}[1]{\textbf{#1}}
\newcommand{\para}[1]{\par\medskip\noindent\textbf{#1}\enspace}
\hypersetup{pdftitle={Odd coefficients at powers of two for A = x + (x A²)': a proof of a conjecture in OEIS A397588},pdfauthor={Dev D. Goyal},pdfsubject={OEIS A397588}}
\title{Odd coefficients at powers of two for $A=x+(xA^2)'$\\[6pt]\large A proof of a conjecture in OEIS A397588}
\author{Dev D. Goyal}
\date{October 2026}
\begin{document}
\maketitle
\begin{abstract}
The sequence A397588 of the On-Line Encyclopedia of Integer Sequences (OEIS), contributed by Paul D.~Hanna, consists of
the coefficients of the power series $A(x)=x+3x^2+24x^3+285x^4+\cdots$ determined by the differential equation
$A=x+(x\,A^2)'$, where the prime denotes the derivative. Hanna conjectured that $a(n)$ is odd exactly when $n$ is a power
of $2$. The entry also states that $a(n)$ is divisible by $3$ for every $n>1$. We prove both statements. Modulo $2$ the
derivative term simplifies, and the equation becomes the Catalan equation $A=x+A^2$, whose solution modulo $2$ is
$x+x^2+x^4+x^8+\cdots$. Modulo $3$, the series $x$ itself is a solution. The results have also been checked by computer
and proved formally in the Lean proof assistant.
\end{abstract}
\begin{center}\begin{minipage}{0.88\linewidth}\small
\textbf{How this paper was produced.} This paper reports the result of an autonomous AI research effort.
ChatGPT (OpenAI), using the GPT-6 model with Pro reasoning and working from prompts by Dev~D.~Goyal, found the proof, wrote its Lean~4 formalization, and wrote a first draft of the paper.
Dev~D.~Goyal chose the problems and asked for checks and revisions, but did not supply proof ideas.
Claude Opus~5.5, used with Claude Code, then checked the work independently: it rebuilt the formal proofs from
source and confirmed that they rely on no unproved assumptions, recomputed the sequence from its definition and tested
every claim against the published data, and read every step of the proofs. It also rewrote this paper for clarity.
Dev~D.~Goyal has verified the results and is responsible for them. The paper has not been peer reviewed.
The section ``About this paper'' at the end gives the details.
\end{minipage}\end{center}

\section{Introduction}

\subsection{The sequence}\label{sec:sequence}
The \emph{On-Line Encyclopedia of Integer Sequences} (OEIS)~\cite{oeisdb} is a large public database of integer
sequences, used by mathematicians and scientists to identify sequences and to record what is known about them. Many
entries also record \emph{conjectures}: patterns that the contributor observed in the data but could not prove.

This paper is about entry A397588~\cite{oeis}, contributed by Paul~D.~Hanna in July 2026. The entry reads, in part:
\begin{quote}\small\ttfamily\raggedright G.f. A(x) satisfies A(x) = x + (x*A(x)\^{}2)'.\par\smallskip Conjecture: a(n) is odd iff n is a power of 2 for n >= 1.\par\smallskip a(n) divisible by 3 for n > 1.\par\smallskip a(1) = 1; a(n) = (n+1) * Sum\_\{k=1..n-1\} a(k)*a(n-k). - Seiichi Manyama, Jul 03 2026\par\smallskip\hfill\textrm{--- OEIS A397588, Paul D. Hanna, July 2026}\end{quote}

Here ``G.f.''\ stands for \emph{generating function}: the terms $a(1),a(2),a(3),\dots$ of the sequence are stored as the
coefficients of the power series $A(x)=a(1)\,x+a(2)\,x^2+a(3)\,x^3+\cdots$. The prime $'$ denotes the derivative.
For power series no limits are needed: the derivative is defined term by term by the familiar rule for powers,
\[
\bigl(u_0+u_1x+u_2x^2+u_3x^3+\cdots\bigr)'=u_1+2u_2x+3u_3x^2+\cdots ,
\]
and it obeys the usual rules, such as the product rule (Section~\ref{sec:background}). So the defining equation is
\begin{equation}\label{eq:def}
A=x+\bigl(x\,A^2\bigr)' .
\end{equation}
Comparing coefficients gives the last formula of the quotation, due to Seiichi~Manyama: $a(1)=1$ and
\[
a(n)=(n+1)\sum_{k=1}^{n-1}a(k)\,a(n-k)\qquad(n\ge2).
\]
Indeed, if $s(n)=\sum_{k=1}^{n-1}a(k)\,a(n-k)$ is the coefficient of $x^n$ in $A^2$, then $x\,A^2=\sum_ns(n)\,x^{n+1}$,
and its derivative is $\sum_n(n+1)\,s(n)\,x^n$. The first terms, starting at $n=1$, are
1, 3, 24, 285, 4284, 75978, 1530720, 34237485, 837481140, \dots

\subsection{The question}
Hanna's conjecture says that $a(n)$ is odd exactly when $n$ is one of $1,2,4,8,16,\dots$. The entry also states, without
calling it a conjecture, that $a(n)$ is divisible by $3$ for every $n>1$. The table shows the remainders modulo $2$ and
modulo $3$ of the first terms. The powers of $2$ are printed in bold.
\begin{center}\small\setlength{\tabcolsep}{3.2pt}
\begin{tabular}{>{\raggedright\arraybackslash}p{6.4em}>{\centering\arraybackslash}p{1.15em}>{\centering\arraybackslash}p{1.15em}>{\centering\arraybackslash}p{1.15em}>{\centering\arraybackslash}p{1.15em}>{\centering\arraybackslash}p{1.15em}>{\centering\arraybackslash}p{1.15em}>{\centering\arraybackslash}p{1.15em}>{\centering\arraybackslash}p{1.15em}>{\centering\arraybackslash}p{1.15em}>{\centering\arraybackslash}p{1.15em}>{\centering\arraybackslash}p{1.15em}>{\centering\arraybackslash}p{1.15em}>{\centering\arraybackslash}p{1.15em}>{\centering\arraybackslash}p{1.15em}>{\centering\arraybackslash}p{1.15em}>{\centering\arraybackslash}p{1.15em}>{\centering\arraybackslash}p{1.15em}>{\centering\arraybackslash}p{1.15em}>{\centering\arraybackslash}p{1.15em}>{\centering\arraybackslash}p{1.15em}}
\toprule
$n$ & \hl{1} & \hl{2} & 3 & \hl{4} & 5 & 6 & 7 & \hl{8} & 9 & 10 & 11 & 12 & 13 & 14 & 15 & \hl{16} & 17 & 18 & 19 & 20 \\ \midrule
$a(n)\bmod 2$ & \hl{1} & \hl{1} & 0 & \hl{1} & 0 & 0 & 0 & \hl{1} & 0 & 0 & 0 & 0 & 0 & 0 & 0 & \hl{1} & 0 & 0 & 0 & 0 \\
$a(n)\bmod 3$ & \hl{1} & \hl{0} & 0 & \hl{0} & 0 & 0 & 0 & \hl{0} & 0 & 0 & 0 & 0 & 0 & 0 & 0 & \hl{0} & 0 & 0 & 0 & 0 \\
\bottomrule
\end{tabular}
\\[4pt]
\begin{tabular}{>{\raggedright\arraybackslash}p{6.4em}>{\centering\arraybackslash}p{1.15em}>{\centering\arraybackslash}p{1.15em}>{\centering\arraybackslash}p{1.15em}>{\centering\arraybackslash}p{1.15em}>{\centering\arraybackslash}p{1.15em}>{\centering\arraybackslash}p{1.15em}>{\centering\arraybackslash}p{1.15em}>{\centering\arraybackslash}p{1.15em}>{\centering\arraybackslash}p{1.15em}>{\centering\arraybackslash}p{1.15em}>{\centering\arraybackslash}p{1.15em}>{\centering\arraybackslash}p{1.15em}>{\centering\arraybackslash}p{1.15em}>{\centering\arraybackslash}p{1.15em}>{\centering\arraybackslash}p{1.15em}>{\centering\arraybackslash}p{1.15em}>{\centering\arraybackslash}p{1.15em}>{\centering\arraybackslash}p{1.15em}>{\centering\arraybackslash}p{1.15em}>{\centering\arraybackslash}p{1.15em}}
\toprule
$n$ & 21 & 22 & 23 & 24 & 25 & 26 & 27 & 28 & 29 & 30 & 31 & \hl{32} & 33 & 34 & 35 & 36 & 37 & 38 & 39 & 40 \\ \midrule
$a(n)\bmod 2$ & 0 & 0 & 0 & 0 & 0 & 0 & 0 & 0 & 0 & 0 & 0 & \hl{1} & 0 & 0 & 0 & 0 & 0 & 0 & 0 & 0 \\
$a(n)\bmod 3$ & 0 & 0 & 0 & 0 & 0 & 0 & 0 & 0 & 0 & 0 & 0 & \hl{0} & 0 & 0 & 0 & 0 & 0 & 0 & 0 & 0 \\
\bottomrule
\end{tabular}
\end{center}

\subsection{What this paper proves}
\textbf{Main result} (Theorem~\ref{thm:main}). \emph{Modulo $2$, the series $A(x)$ equals $x+x^2+x^4+x^8+\cdots$, and
modulo $3$ it equals $x$.} So $a(n)$ is odd if and only if $n$ is a power of $2$, which is Hanna's conjecture, and $a(n)$
is divisible by $3$ for every $n>1$ (Corollary~\ref{cor:conj}).

The proof modulo $2$ rests on the \emph{Catalan numbers} $1,1,2,5,14,42,\dots$ (OEIS A000108), one of the most common
sequences in combinatorics. Their series $C(x)=x+x^2+2x^3+5x^4+\cdots$ is the only series without constant term that
satisfies $C=x+C^2$, and its coefficients are odd exactly at the powers of $2$ (Lemma~\ref{lem:catalan}). Modulo $2$,
equation~\eqref{eq:def} turns into this Catalan equation.

\subsection{What was known}
The entry records several equivalent forms of the equation, the recurrence above (Seiichi~Manyama), and the growth rate
of $a(n)$ (Vaclav~Kotesovec). The divisibility by $3$ is stated in the entry; its proof is a one-line observation, and it
may well have been known to the contributor, so we include it for completeness rather than as a new result. The parity of
the Catalan numbers is classical; see for example~\cite{deutschsagan}. We are not aware of an earlier proof of the parity
conjecture.

\subsection{Why it is interesting}
Equations that combine an unknown power series with its derivative are \emph{differential equations}; they appear
throughout mathematics and its applications, and the OEIS contains many sequences defined by them. This example shows a
special feature of derivatives in modular arithmetic. The derivative multiplies the coefficient of $x^n$ by $n$, so modulo
a prime $p$ it kills every power $x^{n}$ with $n$ divisible by $p$. Both results of this paper come from this effect: modulo
$2$ one term of the equation disappears, and modulo $3$ the derivative of $x^3$ disappears.

\subsection{How to read this paper}
Section~\ref{sec:background} explains the notions we need: power series and their derivatives, remainders, and the
Catalan numbers modulo $2$. Section~\ref{sec:results} states the results precisely, and Section~\ref{sec:proofs} proves
them. Section~\ref{sec:checks} describes how the results were checked, by computer and by a formal proof in Lean.
Section~\ref{sec:discussion} discusses the scope of the result and higher powers of $3$.
\section{Background}\label{sec:background}
This section collects the definitions and facts used later. Readers who know formal power series and congruences can skip it and return when needed.
\subsection{Formal power series}
A \emph{formal power series} is an expression
\[
U(x)=u_0+u_1x+u_2x^2+u_3x^3+\cdots ,
\]
determined by its sequence of coefficients $u_0,u_1,u_2,\dots$. No number is ever substituted for $x$, so questions
of convergence do not arise. We write $\coeff{n}{U}=u_n$ for the coefficient of $x^n$ in $U$; two series are equal
when all their coefficients are equal.

Series are added coefficient by coefficient and multiplied like polynomials:
\[
\coeff{n}{(UV)}=u_0v_n+u_1v_{n-1}+\cdots+u_nv_0 .
\]
Each coefficient of a product is a finite sum. For example, $(1-x)(1+x+x^2+x^3+\cdots)=1$, so $1+x+x^2+\cdots$ is
the \emph{multiplicative inverse} of $1-x$, written $\frac{1}{1-x}$. More generally, every series $D$ with constant term
$1$ has a unique multiplicative inverse $\frac1D$, whose coefficients are determined one at a time by $D\cdot\frac1D=1$
using only addition, subtraction and multiplication. Every fraction in this paper has a denominator of this kind, so it
denotes a power series.

We write $U=O(x^d)$ when all coefficients of $U$ below degree $d$ are zero. So
$U-V=O(x^d)$ means that $U$ and $V$ agree up to degree $d-1$. A series of the form $x+u_2x^2+u_3x^3+\cdots$, with
constant term $0$ and coefficient of $x$ equal to $1$, is called \emph{normalized}.

\subsection{The derivative of a power series}
The \emph{formal derivative} of a power series $U=\sum_nu_nx^n$ is the series
\[
U'=\sum_{n\ge1}n\,u_n\,x^{n-1}=u_1+2u_2x+3u_3x^2+\cdots .
\]
It is defined by this formula alone, with no limits, so it makes sense for coefficients in any commutative ring. It
satisfies the usual rules: $(U+V)'=U'+V'$, and the \emph{product rule} $(UV)'=U'V+UV'$. For the product rule it is enough
to check single powers, where it says $(a+b)\,x^{a+b-1}=a\,x^{a-1}\cdot x^b+x^a\cdot b\,x^{b-1}$; the general case follows
by adding up. In particular $(U^2)'=2\,U\,U'$.

\subsection{Remainders and congruences}
For an integer $m\ge2$ we write $a\equiv b\pmod m$ (``$a$ is congruent to $b$ modulo $m$'') when $a-b$ is divisible by
$m$, that is, when $a$ and $b$ leave the same remainder after division by $m$. For example, $17\equiv5\pmod 4$, since both leave remainder $1$, and $-3\equiv1\pmod4$. The
remainder of $a$ on division by $m$ is called its \emph{residue} modulo $m$; it is one of $0,1,\dots,m-1$. Saying that
$a$ is odd is the same as saying $a\equiv1\pmod2$.

Addition and multiplication are compatible with congruences, so the residues $0,1,\dots,m-1$, added and multiplied
modulo $m$, form a ring, written $\Z/m\Z$. When $m=p$ is a prime we also write $\F_p$. In $\F_2=\{0,1\}$ we have $1+1=0$,
so addition and subtraction coincide.

For power series with integer coefficients, $U\equiv V\pmod m$ means that \emph{every} coefficient of $U-V$ is
divisible by $m$. Equivalently, the two series become equal when every coefficient is replaced by its residue; we call
this \emph{reducing} the series modulo $m$. The key fact is the following.

\emph{Reduction modulo $m$ is compatible with all the operations above.} To reduce a sum, a product, a substitution, or
an inverse, we may reduce the ingredients first and then compute in $\Z/m\Z$. This holds because every
coefficient of the result is obtained from finitely many coefficients of the ingredients by adding, subtracting and
multiplying, and these operations are compatible with remainders.

\para{Rings.} Several lemmas below hold over every \emph{commutative ring}, that is, every set with an addition,
subtraction and multiplication satisfying the usual rules of arithmetic (for example $\Z$, $\Q$ and $\Z/m\Z$). No
division is used; this matters because in $\Z/m\Z$ division is not always possible (for instance, $2$ has no inverse in
$\Z/4\Z$).

\subsection{The method}
Almost every proof in this paper follows the same three steps.
\begin{enumerate}
\item \emph{Uniqueness.} Show that the defining equation determines the coefficients one at a time, using only
addition, subtraction and multiplication. Then the equation has exactly one solution of the required shape, in every
commutative ring, including $\Z/m\Z$.
\item \emph{Reduction.} Reducing the integer solution modulo $m$ gives a solution of the reduced equation,
because reduction is compatible with all operations.
\item \emph{A good guess.} Find a simple series that satisfies the reduced equation and check that it does. By uniqueness,
this simple series \emph{is} the reduction of the sequence, and its coefficients can simply be read off.
\end{enumerate}
The creative part is step~3: finding a simple series that works. Step~1 sometimes needs care, when the equation seems to
require a division.
For a first example, take the Catalan series $C=x+C^2$ from above. Its coefficients $1,1,2,5,14,42,\dots$ are
determined one at a time by this equation. Modulo $2$, the series $S=x+x^2+x^4+x^8+\cdots$ satisfies the same
equation: in $\F_2$, $S^2=x^2+x^4+x^8+\cdots$ (see Lemma~\ref{lem:frobenius}), so $x+S^2=S$. Hence $C\equiv S\pmod2$,
and the Catalan coefficient $\coeff{n}{C}$ is odd exactly when $n$ is a power of $2$.

\subsection{Squaring modulo \texorpdfstring{$2$}{2}}
\begin{lemma}[Frobenius identity]\label{lem:frobenius}
Let $p$ be a prime. For every power series $U$ with coefficients in $\F_p$, we have $U(x)^p=U(x^p)$.
\end{lemma}

\begin{proof}
In $\F_p$, the binomial theorem simplifies to $(a+b)^p=a^p+b^p$, because every binomial coefficient
$\binom pi$ with $0<i<p$ is divisible by $p$. Also $c^p=c$ for every $c\in\F_p$ (Fermat's little theorem). So
$\bigl(\sum_nu_nx^n\bigr)^p=\sum_nu_n^px^{pn}=\sum_nu_nx^{pn}=U(x^p)$. Only finitely many terms matter for any given
coefficient, so this is a polynomial identity.
\end{proof}

For $p=2$ this says that squaring a series modulo $2$ just doubles every exponent: for instance,
$(x+x^3)^2\equiv x^2+x^6\pmod2$. In particular, the square of any series has, modulo $2$, no terms of odd degree.

\subsection{The Catalan numbers modulo 2}
\begin{lemma}\label{lem:catalan}
In every commutative ring the equation $U=x+U^2$ has exactly one solution with $U(0)=0$. Over the integers it is the
Catalan series $C(x)=x+x^2+2x^3+5x^4+14x^5+\cdots$, and
\[
C(x)\equiv x+x^2+x^4+x^8+\cdots=\sum_{k\ge0}x^{2^k}\pmod 2 .
\]
So $\coeff{n}{C}$ is odd exactly when $n$ is a power of $2$.
\end{lemma}

\begin{proof}
The coefficient of $x^n$ in $U^2$ involves only $u_1,\dots,u_{n-1}$, so the equation determines $u_n$ one at a time:
$u_1=1$ and $u_n=\sum_{i=1}^{n-1}u_iu_{n-i}$. Over $\F_2$, the series $S=\sum_kx^{2^k}$ satisfies
$x+S^2=x+\sum_kx^{2^{k+1}}=S$ by Lemma~\ref{lem:frobenius}, so by uniqueness $C\equiv S\pmod2$.
\end{proof}
The parity of the Catalan numbers is classical; see for example~\cite{deutschsagan}. (Here the Catalan number
$\frac1{n}\binom{2n-2}{n-1}$ appears as the coefficient of $x^n$, so ``$n$ a power of $2$'' corresponds to the familiar
statement that the $k$-th Catalan number is odd exactly when $k+1$ is a power of $2$.)

\section{Results}\label{sec:results}

\begin{theorem}\label{thm:main}
\begin{enumerate}[label=(\alph*)]
\item Equation~\eqref{eq:def} has exactly one normalized solution $A$ with integer coefficients (and, more generally,
exactly one with coefficients in any commutative ring). The sequence A397588 is $a(n)=\coeff{n}{A}$.
\item Modulo $2$,
\[
A(x)\equiv x+x^2+x^4+x^8+\cdots=\sum_{k\ge0}x^{2^k}\pmod 2 .
\]
\item Modulo $3$, $A(x)\equiv x\pmod3$.
\end{enumerate}
\end{theorem}

\begin{corollary}\label{cor:conj}
\begin{enumerate}[label=(\roman*)]
\item (Hanna's conjecture.) For every $n\ge1$, $a(n)$ is odd if and only if $n$ is a power of $2$.
\item For every $n>1$, $a(n)$ is divisible by $3$.
\end{enumerate}
\end{corollary}

\begin{proof}
Read off the coefficients in Theorem~\ref{thm:main}(b) and (c).
\end{proof}

\section{Proofs}\label{sec:proofs}

\para{The idea in brief.} The coefficient of $x^n$ in $(xA^2)'$ involves only earlier coefficients of $A$, so the
equation has exactly one solution. By the product rule, $(xA^2)'=A^2+2x\,A\,A'$. Modulo $2$ the second term vanishes, and
the equation becomes the Catalan equation $A=x+A^2$. Modulo $3$ we simply check that $A=x$ is a solution: then
$(x\cdot x^2)'=3x^2$, which is $0$ modulo $3$.

\begin{proof}[Proof of Theorem~\ref{thm:main}]
(a) As computed in the introduction, the coefficient of $x^n$ in $(xA^2)'$ is $(n+1)\sum_{k=1}^{n-1}a(k)\,a(n-k)$, which
involves only $a(1),\dots,a(n-1)$. So \eqref{eq:def} determines the coefficients one at a time: $a(0)=0$, $a(1)=1$, and
$a(n)$ for $n\ge2$ from the earlier ones. Only additions and multiplications are used, so this works in every commutative
ring, and reducing the integer solution modulo $m$ gives the solution over $\Z/m\Z$.

(b) By the product rule, $(xA^2)'=A^2+x\,(A^2)'=A^2+2x\,A\,A'$. Modulo $2$ the last term vanishes, so the reduction of
$A$ modulo $2$ satisfies $A=x+A^2$ over $\F_2$. By Lemma~\ref{lem:catalan}, the only normalized solution of this
equation over $\F_2$ is $\sum_{k\ge0}x^{2^k}$.

(c) The series $x$ is normalized, and over $\Z/3\Z$ it satisfies \eqref{eq:def}: $x+(x\cdot x^2)'=x+3x^2=x$. By the
uniqueness in part (a) for the ring $\Z/3\Z$, the reduction of $A$ modulo $3$ is $x$.
\end{proof}

\section{Checking the results}\label{sec:checks}

A mathematical proof can contain mistakes, so the results were checked in two independent ways.

\para{Numerical check.}
We computed $A$ by repeating the substitution $A\mapsto x+(xA^2)'$ with exact integer arithmetic, starting from
$A=x$. This reproduces all 20 terms printed in the OEIS entry. The parity rule of Corollary~\ref{cor:conj} holds for
every $n<299$, and $3$ divides $a(n)$ for every $n$ with $1<n<299$.
The computation was done with the PARI/GP computer algebra system, using code that was written separately from the
formal proofs (file \texttt{verification/independent\_checks.gp} in the accompanying repository~\cite{repo}).
A computation can only test finitely many cases, so it is a check on the definitions and on the indexing, not a proof.

\para{Formal proof in Lean.}
Lean~\cite{lean} is a \emph{proof assistant}: a program in which mathematical statements and their proofs are written
in a precise formal language, and every step is checked mechanically down to the basic rules of logic. All results of
this paper are proved in Lean~4 (version 4.19.0) on top of Mathlib~\cite{mathlib}, Lean's library of formalized
mathematics. The statements are collected in the file \texttt{lean/Entries/A397588.lean} of the repository,
in the namespace \texttt{OEIS.A397588}. There the derivative is Mathlib's formal derivative \texttt{derivative}, and the proofs are in the module \texttt{lean/Families/DifferentialParity.lean}.
The table lists where each result of this paper is proved.
\begin{center}\small\renewcommand{\arraystretch}{1.25}
\begin{tabular}{>{\raggedright\arraybackslash}p{0.52\linewidth}>{\raggedright\arraybackslash}p{0.4\linewidth}}
\toprule
Statement in this paper & Lean theorem(s)\\
\midrule
Theorem~\ref{thm:main}(a) for integer coefficients; $A$ is normalized and satisfies \eqref{eq:def} & \texttt{differential\_\allowbreak{}existence}, \texttt{differential\_\allowbreak{}definition}, \texttt{differential\_\allowbreak{}source\_\allowbreak{}unique} \\
Corollary~\ref{cor:conj}(i): $a(n)$ is odd if and only if $n$ is a power of $2$ & \texttt{differential\_\allowbreak{}parity} \\
Theorem~\ref{thm:main}(c) and Corollary~\ref{cor:conj}(ii) & \texttt{differential\_\allowbreak{}mod\_\allowbreak{}three}, \texttt{differential\_\allowbreak{}divisible\_\allowbreak{}by\_\allowbreak{}three} \\
\bottomrule
\end{tabular}
\end{center}
The Lean project was rebuilt from source for this paper, and Lean was asked to list the basic assumptions (axioms)
that each of these theorems ultimately depends on. Every one depends only on the three standard axioms of Lean's logic
(\texttt{propext}, \texttt{Classical.choice}, \texttt{Quot.sound}); no step is left unproved. A formal proof certifies
exactly the statement that Lean checks. Whether that statement says what the OEIS entry means is a separate question,
so the Lean definitions were also compared, side by side, with the definition in the OEIS entry
(Section~\ref{sec:sequence}).

\section{Discussion}\label{sec:discussion}

\para{Scope.} Theorem~\ref{thm:main} determines the remainders of every coefficient modulo $2$ and modulo $3$, so it
settles the conjecture of the entry and proves its divisibility statement. It says nothing about the size of $a(n)$.

\para{Higher powers of $3$.} The recurrence gives more. First, $a(n)=(n+1)\,s(n)$ with an integer $s(n)$, so $n+1$
divides $a(n)$ for every $n\ge2$. Second, $9$ divides $a(n)$ for every $n\ge5$. To see this, let $n\ge3$ and split off the
two terms with $k=1$ and $k=n-1$ in $s(n)$:
\[
s(n)=2\,a(n-1)+\sum_{k=2}^{n-2}a(k)\,a(n-k).
\]
Each remaining product is a product of two multiples of $3$, so it is divisible by $9$. Hence $s(n)\equiv2\,a(n-1)
\pmod9$, and if $9$ divides $a(n-1)$, then $9$ divides $s(n)$ and so $a(n)$. Since $a(5)=4284=9\cdot476$, induction
gives $9\mid a(n)$ for all $n\ge5$. A computer calculation suggests that this continues: every fixed power of $3$ seems to
divide $a(n)$ for all large enough $n$. For example, $3^{12}$ divides $a(n)$ for every $n$ with $35\le n<299$. We have not
tried to prove this in general. (These remarks are not part of the Lean formalization.)

\section*{About this paper}\label{sec:about}
\para{How the work was done.}
This paper is one of 26 papers on conjectures from the OEIS that came out of an autonomous AI research effort directed
by Dev~D.~Goyal~\cite{repo}. The roles were as follows.
\begin{itemize}
\item \emph{Dev~D.~Goyal} started the project, chose which problems to pursue, asked for proofs, formal verification and
revisions, and decided what to publish. Dev~D.~Goyal did not supply proof strategies or intermediate arguments.
\item \emph{ChatGPT} (OpenAI), using the GPT-6 model with Pro reasoning, worked autonomously from those prompts. It
searched the OEIS for suitable conjectures, found the proofs, wrote the Lean~4 formalization, and wrote a first draft
of this paper.
\item \emph{Claude Opus~5.5}, used with Claude Code, checked the work independently and then rewrote this paper. It
rebuilt the Lean project from source and confirmed that the theorems listed in Section~\ref{sec:checks} rely only on
Lean's standard axioms; compared the Lean definitions with the OEIS definition; recomputed the sequence with separately
written PARI/GP code, compared it with the terms published in the OEIS, and tested every claim of the paper on many
terms; and read every step of the proofs. It then rewrote the text to make it self-contained and readable. The
theorems and the proof strategies are ChatGPT's. In rewriting, Claude Opus~5.5 expanded the explanations and, in
places, wrote out an argument in more detail or by a more elementary route. The Discussion section was added by Claude Opus~5.5;
the further observations in it were checked by computer and are not part of the Lean formalization.
\end{itemize}
Dev~D.~Goyal has verified the results and is responsible for this paper.

\para{Review status.} This paper has not been peer reviewed, and it has not been submitted to a journal or to the OEIS.

\para{Priority.} An OEIS entry that labels a statement as a conjecture shows only that its contributor did not know
a proof; it does not show that no proof existed. During the research only limited searches for earlier work were
made. No claim of priority is made.

\para{Corrections.} Please report any error in the issue tracker of the repository~\cite{repo}, where the Lean sources and the numerical checks are also
available.

\section*{Acknowledgements}
The sequence, the conjecture and the divisibility statement are due to Paul~D.~Hanna~\cite{oeis}. The recurrence in the
entry was contributed by Seiichi~Manyama, and the asymptotic formula by Vaclav~Kotesovec.

\begin{thebibliography}{99}\small
\raggedright
\bibitem{oeis} P.~D.~Hanna, entry A397588 (G.f.\ $A(x)$ satisfies $A(x)=x+(xA(x)^2)'$), \emph{The On-Line Encyclopedia of Integer Sequences}, July 2026. \url{https://oeis.org/A397588}. Text as of October 2026.
\bibitem{deutschsagan} E.~Deutsch and B.~E.~Sagan, Congruences for Catalan and Motzkin numbers and related sequences, \emph{Journal of Number Theory} 117 (2006), 191--215.
\bibitem{oeisdb} N.~J.~A.~Sloane and The OEIS Foundation Inc., \emph{The On-Line Encyclopedia of Integer Sequences}. \url{https://oeis.org}.
\bibitem{stanley} R.~P.~Stanley, \emph{Enumerative Combinatorics}, Volume~2, Cambridge University Press, 1999. (Chapter~5 treats formal power series, composition and compositional inverses.)
\bibitem{niven} I.~Niven, Formal power series, \emph{American Mathematical Monthly} 76 (1969), 871--889.
\bibitem{lean} L.~de~Moura and S.~Ullrich, The Lean~4 theorem prover and programming language, in \emph{Automated Deduction -- CADE 28}, Lecture Notes in Computer Science 12699, Springer, 2021, 625--635.
\bibitem{mathlib} The mathlib Community, The Lean mathematical library, in \emph{Proceedings of the 9th ACM SIGPLAN International Conference on Certified Programs and Proofs (CPP 2020)}, 367--381. Version used: commit \texttt{c44e0c8ee63c}.
\bibitem{repo} D.~D.~Goyal, \emph{Proofs and corrections of OEIS conjectures, with Lean~4 formalizations}, repository, 2026. \url{https://github.com/dvgyl/oeis-a396843-proof}.
\end{thebibliography}
\end{document}
